\chapter{Selection}

\begin{question}[Selection Problem]
  Given an totally ordered array $S$ of $n$ distinct numbers and an integer $k$, find the $k$-th smallest number in $S$.
\end{question}

When $k = 1$, this is the minimum problem, which can be solved in $O(n)$ time by scanning through the array. Can we reach $O(n)$ time for general $k$?

\section{Randomized Selection}

Inspired by QuickSort, we can use a similar partitioning idea to solve the selection problem:

\begin{algorithm}[H]
  \caption{Rank$(e,S)$}
  scan through $S$ to find the rank of $e$ in $S$\;
  \Return the rank\;
\end{algorithm}

\begin{algorithm}[H]
  \caption{QuickSelect$(k,S)$}
  randomly pick a pivot $e$ from $S$\;
  $r \gets$ Rank$(e,S)$\;
  \If{$k = r$}{
    \Return $e$\;
  }
  $S_<\gets \{x \in S \mid x < e\}$\;
  $S_>\gets \{x \in S \mid x > e\}$\;
  \eIf{$k < r$}{
    \Return QuickSelect$(k,S_<)$\;
  }{
    \Return QuickSelect$(k-r, S_>)$\;
  }
\end{algorithm}

\begin{theorem}
  The expected running time of QuickSelect is $O(n)$.
\end{theorem}
\begin{proof}
  Let $Q(k,n)$ be the expected running time of QuickSelect$(k,S)$ where $|S|=n$. We obtain the following recurrence relation:

  \[
    Q(k,n) = n-1 + \frac{1}{n} \left( \sum_{r = k+1}^{n}Q(k,r-1) + \sum_{r=1}^{k-1} Q(k-r,n-r) \right)
  \]

  For simplicity, we consider the worst case over all $k$ (denote it by $Q(n)$):

  \begin{align*}
    Q(n) &= n-1 + \frac{1}{n} \sum_{r=1}^{n} \max\{Q(r-1), Q(n-r)\}\\
      &\leq n-1 + \frac{2}{n} \sum_{r = \frac{n}{2}}^{n-1} Q(r) \\
  \end{align*}

  We will prove by induction that $Q(n) \leq c n$ for some constant $c > 0$:

  \begin{align*}
    Q(n) &\overset{IH}{\leq} n-1 + \frac{2}{n} \sum_{r = \frac{n}{2}}^{n-1} c r \\
         &= n-1 + \frac{2c}{n}\cdot\frac{3n}{4}\cdot \frac{n}{2}\\
         &\leq n-1 + \frac{3cn}{4} \\
         &\leq cn \text{ when } c \geq 4
  \end{align*}

  Therefore, $Q(n) = O(n)$.
\end{proof}

The worst case, however, is $O(n^2)$ when the pivot is always the smallest or largest element. To avoid this, we can use a deterministic method to choose a good pivot.

\section{Deterministic Selection}

We can leave everything the same as QuickSelect, except for the pivot selection. We will use the \textbf{Median of Medians} algorithm to select a good pivot in linear time.

\begin{algorithm}[H]
  \caption{DeterministicSelect$(k,S)$}
  partition $S$ into groups of 5 elements each (with less than 4 leftover)\;
  $S'\gets$ array of medians of each group\;
  $e \gets$ DeterministicSelect$(\lceil |S'|/2 \rceil, S')$\;
 $r \gets$ Rank$(e,S)$\;
  \If{$k = r$}{
    \Return $e$\;
  }
  $S_<\gets \{x \in S \mid x < e\}$\;
  $S_>\gets \{x \in S \mid x > e\}$\;
  \eIf{$k < r$}{
    \Return DeterministicSelect$(k,S_<)$\;
  }{
    \Return DeterministicSelect$(k-r, S_>)$\;
  }
\end{algorithm}

\begin{theorem}
  The running time of DeterministicSelect is $O(n)$.
\end{theorem}
\begin{proof}
  We first analyze the range of $e$'s rank. 

  Since $e$ is the median of medians, at least half of the groups have their medians less than $e$. Each such group contributes at least 3 elements less than $e$ (the median and two smaller elements). Therefore, the number of elements less than $e$ is at least $3n/10$. 

  Similarly, at least $3n/10$ elements are greater than $e$. Thus, the rank of $e$ is between $3n/10$ and $7n/10$, which also means the sizes of $|S_<|$ and $|S_>|$ are at most $7n/10$.

  Denote $T(n)$ as the running time of DeterministicSelect on an input of size $n$. As finding the median of 5 items takes 8 comparisons, we have the following recurrence relation:

  \[ T(n) = 8n/5 + T(n/5) + n-1 + T(7n/10)\leq 13n/5 + T(n/5) + T(7n/10) \]

  As $1/5 + 7/10 = 9/10 < 1$, the sum of input sizes in the recurrence decays. Therefore, we could expect $O(n)$ to be the dominant term. We will prove by induction that $T(n) \leq c n$ for some constant $c > 0$:

  \begin{align*}
    T(n) &\overset{IH}{\leq} 13n/5 + c n/5 + c \cdot 7n/10 \\
         & = n\left( \frac{13}{5} + \frac{9c}{10} \right) \\
         & \leq c n \text{ when } c \geq 26
  \end{align*}

  Therefore, $T(n) = O(n)$.
\end{proof}

\section{Lower Bound}

Now let's look at the lower bound of the selection problem.
\begin{theorem}
  Any comparison-based algorithm for the selection problem makes at least $\frac{3n}{2}-O(\log n)$ comparisons in the worst case.
\end{theorem}
\begin{proof}
  We use a binary decision tree to model any comparison-based selection algorithm. Each internal node represents a comparison between two elements, and each leaf node represents a possible outcome (i.e., the $k$-th smallest element).

\begin{tikzpicture}[
    node distance=0.2cm and 0.3cm,
    every node/.style={draw, rectangle, rounded corners, align=center, minimum width=2.5cm, minimum height=0.9cm, font=\small},
]

% Nodes
\node (A) {$a\overset{?}{<}b$};
\node (B) [below right=of A] {$a\overset{?}{<}c$};
\node (C) [below left=of A] {$b\overset{?}{<}d$};
\node (D) [below left=of C] {return $e$};
\node (E) [below right=of C] {return $f$};

% Edges
\draw[--, thick] (A) -> (B);
\draw[--, thick] (A) -> (C);
\draw[--, thick] (C) -> (D);
\draw[--, thick] (C) -> (E);

\end{tikzpicture}

The number of comparisons made by the algorithm corresponds to the depth of the leaf node in the decision tree.

Finding the $k$-th smallest element is equivalent to splitting $S$ into 2 sets, $L$ and $U$, where any element in $L$ is less than any element in $U$, and $|L| = k-1$, and then find the smallest element in $U$.

Therefore, the decision tree of Selection must contain subtrees of finding the minimum in $U$ for each possible $L$. Finding the minimum in $n-k+1$ elements requires at least $n-k$ comparisons, so every subtree has depth at least $n-k$, and consequently, at least $2^{n-k}$ leaf nodes.

Then the total number of leaf nodes in the decision tree is at least $\binom{n}{k-1} \cdot 2^{n-k}$, so the number of comparisons (i.e. depth) in the worst case is at least $\log \left( \binom{n}{k-1} \cdot 2^{n-k} \right)$.

This value is maximized when $k = n/2$. According to Stirling's approximation, $\binom{n}{n/2} \approx \frac{2^n}{\Theta(\sqrt{n})}$, so we have:

\[
  \text{number of comparisons} \geq \log \left( \binom{n}{n/2 - 1} \cdot 2^{n - n/2} \right) = \log \left( \frac{2^n}{\Theta(\sqrt{n})} \cdot 2^{n/2} \right) = \frac{3n}{2} - O(\log n)
\]

\end{proof}


